\chapter{The papers read}
\label{chap:C-articles}

This appendix lists the papers of the two samples of the main comparison, with their identifier,
title and venue. The identifier points to the reading file \file{study/codes/<reader>/<id>.json} or
\file{study/codes\_v2/<reader>/<id>.json}, where every decision rests on a quotation and its page.
The PDFs are not in the repository (\cref{sec:absent}). The script \file{study/fetch.py} downloads
them from the addresses recorded in the sample files.

\section{The 2026 papers}
\label{sec:c-2026}

The population consists of the \num{1252} papers of ICLR 2026, ICML 2026 and NeurIPS 2025 whose title
matches the regular expression of the protocol, out of \num{17836} accepted papers
(\cref{chap:03-sondage-methode}). The prefix of the identifier shows where the paper comes from.
\begin{itemize}
  \item \texttt{LR}, \texttt{ML}, \texttt{NE}: the first draw (\texttt{random.Random(2018)}):
        17 papers from ICLR 2026, 13 from ICML 2026 and 12 from NeurIPS 2025
        (\file{study/final\_2026.json}).
  \item \texttt{MLR}, \texttt{NER}: the eight replacement papers from ICML and NeurIPS, drawn by
        \texttt{random.Random(2026)} in the same conference when the drawn paper had no freely
        accessible PDF (\file{replacements\_2026.json}).
  \item \texttt{K}: the queue of enlargement A3, which brings classical RL of 2026 to 50 papers, shuffled by \texttt{random.Random(2026)} and
        read in that order, after the title screen, until 36 classical RL papers had been read. The
        complete queue is in \file{classic\_queue\_2026.json}.
\end{itemize}
The classical RL sample of 2026 in the main comparison combines the 14 \texttt{classic} papers of
the first draw and the 36 of queue A3. \Cref{tab:articles-2026} lists the papers read. These papers were read mostly in their arXiv version, the 2018 papers in the
proceedings. For criterion Q5 alone, the 2026 papers were reread in the proceedings version and the
2018 \enquote{no} decisions in their latest arXiv version (survey deviation 31,
\file{pdf/<id>-proc.pdf} and \file{pdf/<id>-arxiv.pdf}), with the fingerprint of each file in \file{study/pdfs.csv}.

\gentable{articles-2026}

\section{The 2018 papers}
\label{sec:c-2018}

Our reading of 2018 covers 50 RL papers from ICML 2018, NeurIPS 2018 and ICLR 2018, the three
conferences Pineau named. The first 12 form the calibration of the protocol, 6 from ICML and 6 from
NeurIPS (\file{study/calib\_2018.json}). The other 38 come from enlargement A2, which brings the 2018 reading to 50 papers, drawn in the order of
a \texttt{random.Random(2018)} shuffle, with an allocation proportional to the size of each
conference (deviation 12). The sample combines 21 papers from ICML 2018 (prefix \texttt{C1}), 20 from
NeurIPS 2018 (\texttt{C2}) and 9 from ICLR 2018 (\texttt{C3}) (\file{study/calib\_2018\_v2.json}).
The comparison for each criterion counts the papers where the criterion applies, so a denominator
can fall below 50.

\gentable{articles-2018}
